\clearpage
\section{Design}

The first design decision was to avoid treating the three required versions as three different games.
They share the same rules, board, snapshots, and physics semantics.
What changes is only the execution strategy.
This keeps the sequential version useful as a reference and makes comparisons with the two concurrent versions meaningful.

\subsection{A single execution path}

The architecture can be read from left to right:

\begin{center}
\texttt{PooolApplication $\rightarrow$ GameLoop $\rightarrow$ GameModel
$\rightarrow$ Board $\rightarrow$ PhysicsStepper}.
\end{center}

\texttt{PooolApplication} owns the GUI lifecycle and connects the view to a \texttt{GameRuntime}.
In both concurrent versions, \texttt{GameLoop} is the single-writer controller: it accepts queued commands, advances the game, and publishes one immutable snapshot.
\texttt{GameModel} contains match rules, while \texttt{Board} contains the mutable physical state.
The \texttt{PhysicsStepper} interface finally selects how one physical step is computed.

This split gives each class a short answer to the question ``why does it exist?''.
The application assembles components, the loop orders events, the model enforces rules, the board stores the world, and the stepper advances it.

\subsection{Ownership and asynchronous input}

Swing events and the bot can arrive while the simulation is running.
Allowing them to mutate the model directly would make the result depend on timing.
They therefore submit operations to \texttt{CommandMailbox}, a custom monitor with a FIFO queue and completion receipts.
At the beginning of each tick, \texttt{GameLoop} drains the queue on its owner thread.

The ownership rule is deliberately precise: \texttt{GameModel} has one writer.
\texttt{Board} is changed by the physics step invoked by that writer.
Inside a parallel step, workers may update disjoint ball ranges or calculate private collision contributions, but no uncoordinated component can modify the board.
The next snapshot is published only after all phases have joined and committed.

\subsection{One parallel algorithm, two execution policies}

Both concurrent variants use \texttt{ParallelPhysicsEngine}.
The engine contains the complete algorithm and delegates only range execution to a \texttt{RangeScheduler}.
Two scheduler implementations expose the required programming models:

\begin{itemize}
    \item \texttt{PlatformThreadRangeScheduler} uses long-lived
          \texttt{PhysicsWorker} threads and a completion monitor;
    \item \texttt{ExecutorRangeScheduler} submits ranges to a fixed executor
          and joins them through \texttt{Future.get()}.
\end{itemize}

The public \texttt{ThreadedPhysicsEngine} and \texttt{TaskBasedPhysicsEngine} remain as small, explicit facades.
They make the requested variants visible without duplicating collision detection, partitioning, or merge logic.

\subsection{Phased physics and deterministic commit}

A fixed sub-step follows a small number of ordered phases:

\begin{enumerate}
    \item move active balls over disjoint index ranges;
    \item apply hole interactions after the movement barrier;
    \item build one private spatial grid per worker;
    \item merge the grids and sort their cells;
    \item compute collision contributions in private sparse accumulators;
    \item merge deltas, sort contact pairs, and commit the result.
\end{enumerate}

The spatial grid avoids testing every possible pair of balls.
Private grids and accumulators avoid locks in the computational phases.
Sorting cells and contact pairs removes dependence on hash-map iteration order and worker completion order.
Determinism is therefore a property of the merge protocol, not an assumption about thread scheduling.

\subsection{Why the serial sections remain serial}

Command execution, score updates, terminal-state checks, and snapshot publication are short and define the global order of the match.
Parallelizing them would add synchronization without creating useful throughput.
Movement, grid construction, and collision calculation grow with the number of balls, so those are the phases where parallel work pays off.

Small ranges are executed inline, and small merged-delta sets are applied sequentially.
These fast paths avoid paying task and barrier overhead when the work is too small to amortize it.

\subsection{View boundary and shutdown}

The view never renders the live board.
It consumes a \texttt{RuntimeGameSnapshot}, which combines an immutable game snapshot with copied rendering data.
This keeps the Swing Event Dispatch Thread independent of simulation mutations.

Shutdown follows the same ownership rules.
The runner stops producing ticks, the mailbox rejects pending and future commands, a final snapshot is published, and the selected scheduler closes its workers or executor.
No producer is left waiting for a receipt after the runtime has terminated.
